\chapter{强激光场中的经典电子动力学}\label{ch:strong-field-classical}

上一章的含时量子理论适用于任意场强，但在强场电离之后，连续态电子的首要能标常由经典轨道控制。这里不从“三步模型”的经验图像起步，而从带电粒子的最小耦合动力学开始，再依次施加非相对论、偶极、自由连续态和慢包络近似。这样 $U_p$、$3.17U_p$、$2U_p$ 与 $10U_p$ 都能被识别为同一母方程在不同控制条件下的结果。

\section{最小耦合动力学与强场控制参数}

对质量 $m$、电荷 $Q$ 的相对论性点粒子，在给定外电磁场 $(\Phi,\mathbf A)$ 和额外标量势 $V(\mathbf r)$ 中，canonical Hamiltonian 为
\begin{equation}
H(\mathbf r,\mathbf p,t)
=\sqrt{m^2c^4+c^2[\mathbf p-Q\mathbf A(\mathbf r,t)]^2}
+Q\Phi(\mathbf r,t)+V(\mathbf r).
\label{eq:strong-relativistic-hamiltonian}
\end{equation}
机械动量 $\boldsymbol\pi=\mathbf p-Q\mathbf A=\gamma m\mathbf v$ 是可直接进入 Lorentz 力的量。Hamilton 方程等价于
\begin{equation}
\frac{\dd}{\dd t}(\gamma m\mathbf v)
=Q(\mathbf E+\mathbf v\times\mathbf B)-\nabla V.
\label{eq:strong-lorentz-equation}
\end{equation}
对电子取 $Q=-e$、$m=m_e$。

若场近似单色，频率为 $\omega$、峰值电场为 $E_0$，最重要的相对论控制参数是
\begin{equation}
a_0=\frac{eE_0}{m_ec\omega}.
\label{eq:strong-a0}
\end{equation}
$a_0\ll1$ 时 quiver 速度远小于 $c$，可把\eqnrefs{eq:strong-relativistic-hamiltonian}展开到非相对论阶。若原子尺度或电子轨道尺度 $\ell$ 还满足
\begin{equation}
\epsilon_{\rm dip}=k\ell=\frac{\omega\ell}{c}\ll1,
\label{eq:strong-dipole-parameter}
\end{equation}
则在体系范围内可忽略场的空间相位变化，进入 electric-dipole regime。

在这两个控制参数都小时，且电离后的传播阶段暂时忽略母离子势，线偏振场 $\mathbf E(t)=E_0\cos\omega t\,\hat{\mathbf x}$ 给出
\begin{equation}
v(t)
=-\frac{eE_0}{m_e\omega}
[\sin\omega t-\sin\omega t_0],
\label{eq:strong-simpleman-velocity}
\end{equation}
这里额外采用“出生时沿偏振方向机械速度近似为零”的 simple-man 初始条件 $v(t_0)=0$。这一条件不是 Maxwell--Lorentz 动力学的一部分，而是把束缚态隧穿出口与连续态轨道连接起来的模型假设。

\begin{derivation}[从最小耦合 Hamiltonian 到 Lorentz 方程与偶极极限]
相对论最小耦合 Hamiltonian 为
\begin{equation*}
H=\sqrt{m^2c^4+c^2\boldsymbol\pi^2}+Q\Phi+V,
\qquad
\boldsymbol\pi=\mathbf p-Q\mathbf A.
\end{equation*}
由 Hamilton 方程，对位置变量
\begin{equation*}
\dot r_i=\frac{\partial H}{\partial p_i}
=\frac{c^2\pi_i}{\sqrt{m^2c^4+c^2\boldsymbol\pi^2}}
= v_i,
\end{equation*}
即 $\mathbf v=c^2\boldsymbol\pi/H_{\rm kin}$，其中 $H_{\rm kin}=\sqrt{m^2c^4+c^2\boldsymbol\pi^2}$。对正则动量，
\begin{equation*}
\dot p_i=-\frac{\partial H}{\partial r_i}
= -\frac{c^2\pi_j}{H_{\rm kin}}\frac{\partial\pi_j}{\partial r_i}
 -Q\partial_i\Phi-\partial_iV.
\end{equation*}
由 $\partial\pi_j/\partial r_i=-Q\partial_iA_j$ 并代入 $v_j=c^2\pi_j/H_{\rm kin}$，因此
\begin{equation*}
\dot p_i = Qv_j\partial_iA_j - Q\partial_i\Phi - \partial_iV.
\end{equation*}
而机械动量的时间导数为
\begin{equation*}
\dot\pi_i
= \dot p_i - Q\partial_tA_i - Qv_j\partial_jA_i,
\end{equation*}
把 $\dot p_i$ 代回，合并含 $A$ 的两项，
\begin{equation*}
\dot\pi_i
= Q\bigl[-\partial_i\Phi-\partial_tA_i\bigr]
 + Qv_j\bigl(\partial_iA_j-\partial_jA_i\bigr)
 -\partial_iV.
\end{equation*}
由 $\mathbf E=-\nabla\Phi-\partial_t\mathbf A$ 与 $B_i=\epsilon_{ijk}\partial_jA_k$，于是
\begin{equation*}
\dot{\boldsymbol\pi}
= Q\bigl(\mathbf E+\mathbf v\times\mathbf B\bigr)-\nabla V,
\end{equation*}
即\eqnrefs{eq:strong-lorentz-equation}。在非相对论极限 $|\boldsymbol\pi|/(m_ec)=O(a_0)\ll1$ 下，
\begin{equation*}
\gamma = \frac{H_{\rm kin}}{mc^2} = 1+O(a_0^2),
\qquad
\mathbf v = \frac{\boldsymbol\pi}{m}+O(a_0^3 c),
\end{equation*}
故 $\dot{\boldsymbol\pi}=m\dot{\mathbf v}+O(a_0^2)$。再令偶极参数 $k\ell\ll1$，磁场力 $Q\mathbf v\times\mathbf B$ 为 $O(v/c)$ 量级，可忽略。于是
\begin{equation*}
m_e\dot{\mathbf v} = -e\mathbf E(t) - \nabla V.
\end{equation*}
在连续态忽略母离子势 $V$，对线偏振场 $\mathbf E(t)=E_0\cos\omega t\,\hat{\mathbf x}$ 积分，
\begin{equation*}
v_x(t) = -\frac{eE_0}{m_e\omega}\sin\omega t + C.
\end{equation*}
施加 simple-man 初始条件 $v(t_0)=0$，定出 $C=(eE_0/m_e\omega)\sin\omega t_0$，因此
\begin{equation*}
v_x(t) = -\frac{eE_0}{m_e\omega}\bigl[\sin\omega t-\sin\omega t_0\bigr],
\end{equation*}
即\eqnrefs{eq:strong-simpleman-velocity}。
\end{derivation}

\section{质动力势与慢--快变量分离}

完全自由、无漂移电子的 quiver 速度幅值为 $eE_0/(m_e\omega)$，周期平均动能定义为 ponderomotive energy
\begin{equation}
\boxed{
U_p=\left\langle\frac12m_ev_q^2\right\rangle
=\frac{e^2E_0^2}{4m_e\omega^2}.
}
\label{eq:ponderomotive-energy}
\end{equation}
相应 quiver displacement 幅值
\begin{equation}
\alpha_0=\frac{eE_0}{m_e\omega^2}
\label{eq:quiver-amplitude}
\end{equation}
给出另一个空间控制参数。若包络在长度 $L$ 上变化并满足 $\alpha_0/L\ll1$，可把轨道写成慢漂移 $\mathbf R(t)$ 与快振荡 $\boldsymbol\xi(t)$ 的和。

\begin{derivation}[从多尺度展开得到质动力，即\eqnrefs{eq:ponderomotive-force}]
将电子运动分解为快振荡（quiver）和慢漂移（drift），对快运动求精确解，代入运动方程后做周期平均，残余项给出慢运动的有效势。这是多尺度方法的标准范式：引入快时间 $\tau=\omega t$ 与慢时间 $T=t$，把 $\partial_t$ 替换为 $\omega\partial_\tau+\partial_T$，按 $\epsilon=1/(\omega T_{\rm env})\ll1$ 展开。

先讨论分解快慢运动。以一维场 $E(x,t)=E(X)\cos\omega t$ 为例，令 $x=X+\xi$，其中 $X$ 在一个光学周期内近似不变（慢坐标），$\xi$ 是快振荡（quiver）。控制参数 $\alpha_0/L\ll1$ 保证场在 quiver 幅值范围内可线性化。

随后讨论最低阶快运动。在 $X$ 固定的近似下，快运动满足
\begin{equation*}
 m_e\ddot\xi=-eE(X)\cos\omega t,
\end{equation*}
积分两次得到
\begin{equation*}
 \xi=\frac{eE(X)}{m_e\omega^2}\cos\omega t,
\end{equation*}
即 quiver 幅值 $\alpha_0=eE(X)/(m_e\omega^2)$。这里已选择 $\xi$ 的零均值为慢坐标 $X$ 的定义，这一约定把快慢分量正交分离。quiver 速度幅值 $v_q=eE(X)/(m_e\omega)=c\,a_0$，其中 $a_0$ 是\eqnrefs{eq:strong-a0} 的相对论参数；$a_0\ll1$ 保证 $v_q\ll c$，非相对论近似自洽。

再讨论场在慢坐标附近展开。将场在 $X$ 附近 Taylor 展开到一阶：
\begin{equation*}
 E(X+\xi,t)
 =\left[E(X)+\xi E'(X)+O((\xi/L)^2)\right]\cos\omega t.
\end{equation*}
运动方程 $m_e\ddot{x}=-eE(x,t)$ 分解为快慢两部分。快部分已在第二步处理；慢部分取周期平均。二阶项 $O((\xi/L)^2)$ 在 $\alpha_0/L\ll1$ 下属于高阶修正，会在下一阶产生 ponderomotive 势的非线性修正和快慢耦合。

接着讨论周期平均。零阶项 $\langle E(X)\cos\omega t\rangle=0$。一阶项
\begin{equation*}
 \langle m_e\ddot X\rangle
 =-eE'(X)\langle\xi\cos\omega t\rangle
 =-eE'(X)\frac{eE(X)}{m_e\omega^2}\langle\cos^2\omega t\rangle
 =-\frac{e^2E(X)E'(X)}{2m_e\omega^2}
 =-\frac{e^2}{4m_e\omega^2}\frac{\dd E^2}{\dd X}.
\end{equation*}
（$\langle\cos^2\omega t\rangle=1/2$。）这一步的关键是 $\langle\xi\cos\omega t\rangle\ne0$：quiver 与驱动场同相位，所以快振荡与场幅值的乘积在周期平均后留下非零直流分量。若改用圆偏振场，$\langle\xi\cdot E\rangle$ 的几何因子改变但 $U_p\propto E_0^2/(4\omega^2)$ 的结构不变。圆偏振场的 quiver 在偏振面内做圆周运动，$U_p$ 仍正比于 $E_0^2$ 但系数与线偏振差一个因子 2。

进一步讨论有效势。因此到 $O(\alpha_0/L)$，
\begin{equation}\label{eq:ponderomotive-force}
 \mathbf F_p=-\nabla U_p(\mathbf R),
 \qquad
 U_p(\mathbf R)=\frac{e^2E_0^2(\mathbf R)}{4m_e\omega^2}.
\end{equation}
即\eqnrefs{eq:ponderomotive-force}。电子平均被推向场强较弱处（$U_p$ 正比于 $E^2$，力指向 $E$ 减小方向）。若 $\alpha_0/L$ 不小，局域 $U_p$ 不再足以描述慢运动，需保留更高阶修正。高阶修正包括 $U_p$ 的非线性梯度项和快慢耦合项，在相对论强场（$a_0\sim1$）时还需考虑磁漂移修正。

质动力势还是 Keldysh 参数 $\gamma_K=\sqrt{I_p/(2U_p)}$ 的能标：$\gamma_K\gg1$ 对应多光子区，$\gamma_K\ll1$ 对应准静态隧穿区。若空间非均匀尺度不再远大于 quiver 振幅，或相对论参数接近一，则必须保留更高阶梯度、磁漂移和相对论修正。
\end{derivation}

\section{回碰约束与 HHG/ATI 截止}

在偶极、非相对论和自由连续态近似下，把电离相位和第一次回碰相位写成
\[
\phi_0=\omega t_0,
\qquad
\phi_r=\omega t_r.
\]
由\eqnrefs{eq:strong-simpleman-velocity}再积分一次，回到母离子位置的几何条件为
\begin{equation}
F(\phi_0,\phi_r)
=\cos\phi_r-\cos\phi_0
+(\phi_r-\phi_0)\sin\phi_0=0.
\label{eq:strong-return-constraint}
\end{equation}
回碰动能用 $U_p$ 无量纲化后为
\begin{equation}
\frac{K_r}{U_p}
=2(\sin\phi_r-\sin\phi_0)^2.
\label{eq:strong-return-energy}
\end{equation}
在第一次返回支上对\eqnrefs{eq:strong-return-energy}作约束极值，得到
\begin{equation}
K_r^{\max}\simeq3.17U_p,
\qquad
\hbar\omega_{\rm HHG}^{\max}=I_p+3.17U_p.
\label{eq:hhg-cutoff}
\end{equation}

若电子不回碰而直接离开，脉冲前后 $\mathbf A\to0$ 时最终漂移动能满足 $K_{\rm direct}\le2U_p$。若第一次回碰近似为弹性后向散射，则末态边界条件改变为速度反号，得到
\begin{equation}
\frac{K_{\rm rescatt}}{U_p}
=2[2\sin\phi_r-\sin\phi_0]^2,
\qquad
K_{\rm rescatt}^{\max}\simeq10U_p.
\label{eq:ati-cutoffs}
\end{equation}
三个著名截止因此不是独立经验律，而是同一回碰约束配合三种末态条件的极值。

\begin{figure}[tbp]
\centering
\includegraphics[width=.83\textwidth]{strong_field_return.pdf}
\caption{第一返回支上的经典强场截止。$3.17U_p$ 对应回碰动能极值，$10U_p$ 对应回碰后近似反向弹性散射的末态动能极值。}
\label{fig:strong-field-return-energy}
\end{figure}

\begin{derivation}[从回碰约束到三个截止能标]
先讨论回碰约束的几何含义。由\eqnrefs{eq:strong-simpleman-velocity}再积分一次，电子从 $t_0$ 时刻出生、$t_r$ 时刻回到母离子位置的几何条件为\eqnrefs{eq:strong-return-constraint}，即
\[
F(\phi_0,\phi_r)
=\cos\phi_r-\cos\phi_0+(\phi_r-\phi_0)\sin\phi_0=0.
\]
物理上，$\cos\phi_r-\cos\phi_0$ 是 quiver 位移，$(\phi_r-\phi_0)\sin\phi_0$ 是 drift 位移，两者抵消才能回到原点。无量纲回碰动能为
\[
\mathcal K_r(\phi_0,\phi_r)
=2(\sin\phi_r-\sin\phi_0)^2.
\]
第一返回支由 $F=0$ 且 $t_r-t_0$ 最小的支选定。

随后讨论HHG 截止的约束极值。约束极值满足 Lagrange 条件
\[
\nabla\mathcal K_r+\Lambda\nabla F=0,
\qquad F=0.
\]
数值解给出第一返回支的最优相位约为
\[
\phi_0\simeq0.31,
\qquad
\phi_r\simeq4.40,
\]
即电子在激光相位 $\phi_0\simeq0.31$（接近场峰但略偏）时刻隧穿电离，经过约 $4.09$ 弧度的激光周期后回碰母离子。代入得 $\mathcal K_r\simeq3.17$。若回碰时复合到基态，光子能量还需加回束缚能 $I_p$（能量守恒：电子动能 $+$ 光子能量 $=I_p$），得到\eqnrefs{eq:hhg-cutoff}，即 $\hbar\omega_{\rm HHG}^{\max}=I_p+3.17U_p$。$3.17$ 不是普适常数，而是 simple-man 轨道在第一返回支上的特定极值；高返回支给出更低的局部极大，不主导 HHG 谱。

约束极值条件 $\nabla\mathcal K_r+\Lambda\nabla F=0$ 展开为
\begin{equation*}
 \frac{\partial\mathcal K_r}{\partial\phi_0}+\Lambda\frac{\partial F}{\partial\phi_0}=0,
 \qquad
 \frac{\partial\mathcal K_r}{\partial\phi_r}+\Lambda\frac{\partial F}{\partial\phi_r}=0,
\end{equation*}
加上约束 $F=0$ 共三个方程解三个未知数 $\phi_0,\phi_r,\Lambda$。数值解给出第一返回支的唯一极大点；第二返回支的极大值约 $1.8U_p$，远低于第一支，故 HHG 谱截止由第一返回支主导。

再讨论直接 ATI 截止。对不回碰而直接离开的电子，最终 canonical momentum 守恒给出 $p_f=-eA(t_0)$（脉冲消失后 canonical = mechanical），而 $|A|\le E_0/\omega$，故
\[
K_f\le\frac{e^2E_0^2}{2m_e\omega^2}=2U_p.
\]
极值在 $A(t_0)=\pm E_0/\omega$ 即 $\phi_0=\pi/2$ 或 $3\pi/2$（场峰时刻出生）达到。物理上，场峰时刻出生的电子获得最大 quiver 动能，脉冲消失后全部转化为漂移动能。

接着讨论rescattering ATI 截止。若回碰瞬间机械速度反向（弹性后向散射近似），散射后传播到场消失时
\[
p_f=-e[2A(t_r)-A(t_0)],
\]
从而得到\eqnrefs{eq:ati-cutoffs}中的无量纲能量 $\mathcal K_{\rm rescatt}=2[2\sin\phi_r-\sin\phi_0]^2$。继续在同一 $F=0$ 第一返回支上极大化便得到约 $10U_p$，最优相位与 HHG 截止不同。\figref{fig:strong-field-return-energy}给出这两个约束极值的数值轨迹。

以氢原子 \(I_p=13.6\,\si{eV}\) 和 \(800\,\si{nm}\) 激光为例，强度 \(I=10^{14}\,\si{W/cm^2}\) 给出 \(U_p\approx5.97\,\si{eV}\)。因此 HHG 截止约为 \(I_p+3.17U_p\approx32.5\,\si{eV}\)，对应第 \(21\) 阶附近；直接与回碰 ATI 截止分别约为 \(11.9\,\si{eV}\) 和 \(59.7\,\si{eV}\)。此时 \(\gamma_K=\sqrt{I_p/(2U_p)}\approx1.07\)，体系位于多光子与隧穿区的交界，Coulomb 修正和复时间轨道均不可忽略。

三个截止 $3.17U_p$、$2U_p$、$10U_p$ 来自同一回碰约束\eqnrefs{eq:strong-return-constraint}配合不同末态条件；在 SFA 中，这些经典轨道重新表现为复时间鞍点。
\end{derivation}

\section{适用域与通向量子强场理论}

上述 classical simple-man 链条需要同时满足若干独立条件：$a_0\ll1$ 保证非相对论，$k\ell\ll1$ 保证偶极近似，连续态传播中激光力相对 Coulomb 力占主导，且脉冲包络在单个回碰周期内变化不剧烈。母离子势会导致 Coulomb focusing、轨道相位修正和低能结构；有限横向隧穿动量会降低实际回碰概率；当 $a_0$ 接近一时，磁漂移与相对论质量修正都不可忽略。

强场参数常用 Keldysh parameter
\begin{equation}
\gamma_K=\sqrt{\frac{I_p}{2U_p}}
=\frac{\omega\sqrt{2m_eI_p}}{eE_0}
\label{eq:keldysh-parameter-classical}
\end{equation}
比较束缚态隧穿时间尺度与外场周期。$\gamma_K\gg1$ 更接近多光子区，$\gamma_K\ll1$ 更接近准静态隧穿区；它并不单独保证 simple-man 轨道的全部近似成立。下一章将从规范协变的量子传播出发，把本章轨道提升为 Volkov/SFA 的复时间量子轨道，并系统说明这些经典截止为什么会在鞍点极限中重新出现。

\begin{derivation}[ADK 隧穿电离率与 Keldysh 理论的连接]
ADK (Ammosov--Delone--Krainov) 公式给出原子在强场下的隧穿电离率，是 Keldysh 理论在 \(\gamma_K\ll1\) 极限下的半经典结果。分四步建立。

先讨论准静态隧穿图像。在 \(\gamma_K\ll1\)（低频强场），电离可视为准静态隧穿：电子在瞬时电场 \(E\) 下穿越由 Coulomb 势和外场势形成的势垒。势能 \(V(r)=-\frac{Z_{\rm eff}e^2}{r}-eEr\)（沿场方向），势垒高度 \(V_{\rm max}=-2\sqrt{Z_{\rm eff}e^3E}\)（Sommerfeld 参数），势垒宽度 \(r_{\rm tun}=I_p/(eE)\)。WKB 隧穿指数
\begin{equation*}
  \Gamma\sim\exp\!\left[-2\int_0^{r_{\rm tun}}\sqrt{2m_e(I_p-V(r))}\,\dd r\right]\sim\exp\!\left[-\frac{2(2m_e)^{1/2}I_p^{3/2}}{3\hbar eE}\right].
\end{equation*}
定义 Keldysh 指数 \(\kappa=\sqrt{2m_eI_p}/\hbar\) 和场参数 \(F_0=I_p^2/(4Z_{\rm eff}e^3\hbar)\)（原子单位下 \(F_0=E_a/(4Z_{\rm eff})\)），隧穿指数化为 \(-\frac{4\sqrt{2Z_{\rm eff}}}{3}\cdot\frac{F_0}{E}\)。

随后讨论ADK 公式。对量子数为 \((l,m)\) 的态，ADK 电离率
\begin{equation}
  \Gamma_{\rm ADK}=\frac{C_{n^*l}^2}{2\pi n^*}f(l,m)I_p\left(\frac{2F_0}{E}\right)^{2n^*-|m|-1}\exp\!\left[-\frac{2F_0}{3E}\right],
  \label{eq:strong-field-adk}
\end{equation}
其中 \(n^*=Z_{\rm eff}/\sqrt{2I_p}\) 是有效主量子数，\(C_{n^*l}=(2e/n^*)^{n^*}/\Gamma(n^*+l+1)\) 是渐近归一化常数，\(f(l,m)=(2l+1)(l+|m|)!/[2^{|m|}|m|!(l-|m|)!]\)。指数项是 WKB 隧穿，前因子来自 Coulomb 修正和角向投影。

再讨论Keldysh 理论的统一框架。Keldysh 1965 给出任意 \(\gamma_K\) 的电离率
\begin{equation*}
  \Gamma_{\rm Keldysh}\sim\exp\!\left[-\frac{2F_0}{3E}g(\gamma_K)\right],
\end{equation*}
其中 \(g(\gamma)=\frac{3\gamma-\gamma^3}{2}\arcsin\frac{1}{\gamma}-\frac{3}{2}\sqrt{\gamma^2-1}\)（\(\gamma>1\)，多光子区）和 \(g(\gamma)=1-\frac{\gamma^2}{6}\)（\(\gamma<1\)，隧穿区）。极限：\(\gamma\to0\) 时 \(g\to1\)，恢复 ADK 指数 \(-2F_0/(3E)\)；\(\gamma\to\infty\) 时 \(g\sim\frac{3\gamma}{2}\arcsin(1/\gamma)\to\frac{3}{2}\)，指数化为多光子幂律 \(\Gamma\sim E^{2n}\)（\(n=\lceil I_p/\hbar\omega\rceil\)，最低阶多光子过程）。

接着讨论非绝热修正与 Coulomb 效应。ADK 是准静态近似（瞬时场冻结），有限频率修正（非绝热）来自 Keldysh 含时隧穿：\(\Gamma_{\rm nonad}=\Gamma_{\rm ADK}[1+O(\gamma_K^2)]\)。Coulomb 修正（Yudin--Ivanov 2001，Tong--Lin 2005）在低能区显著：\(\Gamma_{\rm Coulomb}=\Gamma_{\rm ADK}\cdot\exp[C(\gamma_K)/\hbar\omega]\)，\(C\) 是 Coulomb 相位修正，在 \(\gamma_K\sim1\) 时可达量级修正。这些修正解释了 ADK 在阈值附近（\(\gamma_K\sim1\)）的系统偏差，是现代强场电离理论（Perelomov--Popov--Terent'ev, PPT）的改进方向。ADK 公式在 \(\gamma_K\ll1\) 且 \(E\ll E_a\)（场远小于原子场）时精度优于 10\%，是强场原子物理的标准工具。

在过渡区 $\gamma_K\sim1$，不能把 ADK 的准静态展开或最低阶多光子幂律单独使用；应保留完整的 Keldysh/PPT 作用量，并按所需精度加入 Coulomb 与非绝热修正。










Corkum 三步模型把电离、连续传播与复合连接起来，并给出第一返回支的 \(3.17\,U_p\) 截止。

先讨论隧穿电离。强场下原子势垒被激光场压低，电子通过隧穿电离进入连续态。电离率由 ADK 公式给出 \(w_{\rm ADK}\sim\exp(-2(2I_p)^{3/2}/3E)\)（\(I_p\) 电离能，\(E\) 场强）。电离时刻 \(t_i\) 由场强决定：电子在 \(E(t_i)=E_0\sin\omega t_i\) 最大时电离概率最高，但截止能来自特定电离相位。

随后讨论真空加速。电离后电子在激光场中经典运动（忽略 Coulomb 回收）：\(\ddot{x}=-E_0\sin\omega t\)，解为 \(v(t)=-\frac{E_0}{\omega}(\cos\omega t-\cos\omega t_i)\)，\(x(t)=\frac{E_0}{\omega^2}(\sin\omega t-\sin\omega t_i-\omega(t-t_i)\cos\omega t_i)\)。电子动能 \(K(t)=\frac{v^2}{2}=\frac{2U_p}{1}(\cos\omega t-\cos\omega t_i)^2\)（\(U_p=E_0^2/4\omega^2\) 有质动力能）。

电子返回原点 \(x(t_r)=0\) 时发生复合，发射光子能量为 \(\hbar\omega_{\rm HHG}=I_p+K(t_r)\)。截止能来自对电离相位 \(\omega t_i\) 最大化 \(K(t_r)\)。返回条件 \(\sin\omega t_r-\sin\omega t_i=(\omega t_r-\omega t_i)\cos\omega t_i\)，数值优化给出最佳电离相位 \(\omega t_i\approx17^\circ\) 和返回相位 \(\omega t_r\approx247^\circ\)，最大动能 \(K_{\max}\approx3.17\,U_p\)。截止能
\begin{equation}
  \hbar\omega_{\rm cutoff}=I_p+3.17\,U_p
  \label{eq:phys-sf-hhg-cutoff}
\end{equation}
是三步模型的标志性预言，与实验精确一致。

接着讨论量子修正与半经典极限。三步模型是半经典近似（电子在连续态用经典轨道），量子修正来自：(1) 隧穿电离的 WKB 相位（给电离率指数前因子）；(2) 传播子的量子扩散（给复合振幅的 \(1/\sqrt{|x''(t_r)|}\) 因子）；(3) Coulomb 回收修正（给截止能附近谱形的对数修正）。Lewenstein 量子理论的半经典极限精确给出三步模型，量子修正系统可计算——这把 HHG 从唯象模型提升为可控近似的严格理论，三步模型的 \(3.17\,U_p\) 截止能是半经典方法的成功典范。
\end{derivation}
